% !TEX root = ../Main.tex
\chapter{数学 \quad 15:00--17:00}

\section{时间分配表}

\begin{center}
\begin{tabular}{@{} >{\bfseries}c l @{}}
    \toprule
    时间段 & 分块 \\
    \midrule
    15:00--15:20(30) & 选择题 \\
    15:20(30)--15:35(50) & 填空题 \\
    15:50--16:20 & 两道大题 \\
    16:20--16:40 & 三道大题 \\
    16:40--16:55 & 五道大题 \\
    16:55--17:00 & 收尾、检查 \\
    \bottomrule
\end{tabular}
\end{center}

\section{注意事项}

\rmkb{
    \begin{enumerate}
        \item 选择和填空部分要\textbf{权衡好做题时间与正确率之间的关系}。考试时还要引入一个选填部分\textbf{"性价比"}高的因素综合考虑。在后期，选填部分要\textbf{留心关注}，应对不同情况形成不同对策。作答时间建议在 \textbf{35--50 分钟}。
        \item 17--19 题平均用时应控制在 \textbf{10 分钟内}；20--22 题平均用时为 \textbf{15 分钟}，其他压轴题应自己要心里有数；两部分\textbf{之内}的题可以灵活变更做题顺序。
        \item 要学会\textbf{全局把握}——争取在\textbf{部分分值}内拿高分，而不是争取在\textbf{卷面分值}中拿高分。
    \end{enumerate}
}
